\documentclass[11pt]{article}
\usepackage{amsmath,amssymb}
\title{Constant flux: a logarithmic spin bound and the exact remaining magnetic determinant}\author{}\date{}
\begin{document}\maketitle
\paragraph{Actual background and full gauge quotient.}
Keep the complete original relative cube, boundary conditions, and exact
nearby Cartan minimum of the complete cube abelian signed incidence calculation:
$a_{(x,\mu)}=\mathbf1_{\mu=1}x_0/R^2$,
$U_e^*(\epsilon)=\exp(i\epsilon a_e\sigma_3)$.
Write $f=R^{-2}$. The complete classical energy
$E_R=(R+1)^2/R^2$ includes its fixed-only part $4/R$.
Use all free links before tree fixing.
In the charged right frame the interior covariant gauge derivative is
\[
 (D_\epsilon\xi)_e=e^{-2i\epsilon a_e}\xi_x-\xi_{x+\hat\mu},
\]
with zero boundary vertex parameters. Let $C_\epsilon$ be the actual
charged Wilson Hessian, and set
\[
 Q_\epsilon=D_\epsilon^*D_\epsilon,\qquad
 V_\epsilon=C_\epsilon+D_\epsilon D_\epsilon^*.
\]
The neutral operators are defined in the same way. The exact quadratic
gauge quotient, with its original unit tree Jacobian, gives
\[
 F_{\rm quad}=\tfrac12\log\det V_{\rm neutral}
 -\log\det Q_{\rm neutral}
 +\log\det V_\epsilon-2\log\det Q_\epsilon+\text{constant}.
\]
This preserves the original raw Gaussian prefactor; it does not identify
that prefactor with the nonlinear Wilson integral.

\paragraph{A literal component decomposition.}
Group free links by their direction $\mu$. Direct combination of the
four ordered Wilson plaquette insertions and the two endpoint gauge
insertions yields
\[
 V_\epsilon=L_\epsilon+P_\epsilon+S_\epsilon.                 \tag{1}
\]
Here $L$ is block diagonal in $\mu$. Each block has the original flat
mixed-relative diagonal and nearest-neighbor hoppings of modulus one,
with the phases inherited from the displayed $D$ and the original Wilson
Hessian. It is a scalar magnetic difference operator on the staggered
edge grid: Neumann in direction $\mu$, Dirichlet in the other three.
Its phases are not replaced by a continuum field. The diagonal correction is
\[
 (P_\epsilon)_{ee}=
 \begin{cases}2(\cos(\epsilon f)-1),&\mu=0,1,\\0,&\mu=2,3.\end{cases}   \tag{2}
\]
Finally $S$ has only off-component entries. A pair of perpendicular
free links sharing a vertex receives one plaquette and one gauge term.
At zero flux they cancel. At nonzero flux the entry is the difference of
two unit phases whose exponents differ by $\epsilon f$ up to sign.
Only the 01 component pair survives for this background. In particular
\[
 S_0=0,\qquad \|S'_0\|\le 4f.                                    \tag{3}
\]
The last bound follows from the at most four adjacent perpendicular
links per row and column. Signs/phase conventions do not affect it.
The same-component diagonal (2) counts the two 01 plaquettes incident
to each free 0- or 1-link. Boundary normal links have one fewer gauge
endpoint; this is already in the original flat Neumann diagonal, not
an omitted boundary correction.

At zero let $G=V_0^{-1}=\operatorname{diag}(G_\mu)$ and $g=Q_0^{-1}$.
Because $L'$ is component diagonal and $S'$ is off component,
\[
 \operatorname{Tr}G S''=0,\qquad
 \operatorname{Tr}(GL'GS')=0.                                    \tag{4}
\]
These are exact finite-matrix identities. Thus
\[
 F_{\rm quad}''=\mathcal M_R-\operatorname{Tr}(GS'GS')
                  +\operatorname{Tr}(GP'')+\mathcal N_R,           \tag{5}
\]
\[
 \mathcal M_R=
 \operatorname{Tr}(GL''-GL'GL')
 -2\operatorname{Tr}(gQ''-gQ'gQ'),\qquad
 \mathcal N_R=-\tfrac12 f^2\operatorname{Tr}(G H_{01}).
\]
Here $H_{01}$ is the complete flat 01 plaquette curl Gram matrix.
The neutral ghost is constant. In particular no magnetic/spin mixed
bubble needs an absolute-value estimate.

\paragraph{An actual logarithmic bound for the spin part.}
The complete relative product spectrum already gives
\[
 \operatorname{Tr}G\le CR^4,\qquad
 \operatorname{Tr}G^2\le CR^4\log(2+R).
\]
Indeed inverse-square eigenvalues are bounded by
$CR^4/|k|^4$, whose four-dimensional shell sum is logarithmic.
Since $S'$ is Hermitian, in an eigenbasis of $G$ the inequality
$g_i g_j\le(g_i^2+g_j^2)/2$ proves
\[
 0\le\operatorname{Tr}(GS'GS')
 \le\|S'\|^2\operatorname{Tr}G^2\le C\log(2+R).                    \tag{6}
\]
No operator-order squaring or commutativity is assumed.
Also $0\le H_{01}\le V_0$, so
\[
 |\mathcal N_R|\le C f^2 R^4\le C,\qquad
 |\operatorname{Tr}(GP'')|\le C f^2\operatorname{Tr}G\le C.         \tag{7}
\]
Consequently the original signed quadratic response satisfies the
proved reduction
\[
 \boxed{\quad |F_{\rm quad}''-\mathcal M_R|
                  \le C[1+\log(2+R)].\quad}                       \tag{8}
\]
For canonical $\epsilon=\zeta/\sqrt B$, divide (5)--(8) by $B$.
As $E_R\ge1$, this is a physical-energy-normalized logarithmic estimate
for every term EXCEPT the explicit scalar magnetic combination
$\mathcal M_R$. It is not a logarithmic bound on the full response.

\paragraph{Independent finite checks and the remaining estimate.}
The finite calculation constructs all
free-link charged Hessian entries, $D,D',D''$, and the complete gauge
quotient independently of the tree inverse. Its values are
\[
\begin{array}{c|r|r|r|r|r}
R&\mathcal M_R&-\operatorname{Tr}(GS'GS')&
\operatorname{Tr}(GP'')&\mathcal N_R&F_{\rm quad}''\\ \hline
2&0&-0.013888889&-0.072916667&-0.036458333&-0.123263889\\
3&0.001407762&-0.060212216&-0.178914364&-0.082304527&-0.320023345
\end{array}
\]
The sum agrees with the direct complete gauge determinant through
$1.2\,10^{-16}$ and with the independent tree calculation.
The mixed boundary scalar magnetic determinants remain exactly in
$\mathcal M_R$. A closed-loop area or finite-boundary magnetic Green
estimate must control their SIGNED second derivatives; their separate
potential-coordinate bounds do not prove (8) with $\mathcal M_R$ removed.
Even that additional estimate would leave the exact original-law
nonlinear residual of the signed gauge quotient and unsummed current calculation(5). No full S94 or scale induction is
claimed here.
\end{document}
